% 本文件由 LaTeX 排版 Agent 根据有效 translation.json 自动生成。
% author=Augustine of Hippo; work_id=1402; title=论摩尼派之道德

\chapter{论摩尼派之道德}

% structure: p0001 source=h1 target=omitted reason=page-title-already-rendered-by-parent

\Emph{由希波的圣奥斯定写于公元388年。}\par

\Emph{包含对这些异端者关于恶之起源与本质之教义的特别驳斥；揭露他们所谓的口、手、胸的象征性习俗；并谴责他们迷信的禁食和不洁的奥秘。最后，提及摩尼派中被揭露的某些罪行。}\par

% structure: p0004 source=h2 target=section reason=relative-heading-rank; base=h2

\section{第一章。——至善即为拥有至高存在者。}

1. 我想，人人都会同意，善恶问题属于道德科学，其中这些术语常用。因此，但愿人们能带着如此清晰的理智完美来从事这些探究，以便能够看见至善——没有什么比它更好或更高，次于它的是处于纯洁与完美状态的理性灵魂。如果这一点被清楚地理解，那么也会显而易见：至善就是被恰当地描述为拥有至高和原初存在的那一位。因为，最高意义上的存在，是指始终如一、处处相似、不可在任何部分被败坏或改变、不受时间限制、在当下与先前状态之间不允许任何变化的东西。这才是真正的存在。因为在\Quote{存在}一词的这个意义上，隐含着一个自足且不变的本质。这样的属性只能归于天主，严格来说没有什么与他相反。因为存在的对立面是非存在。因此，没有与天主相反的本质。但由于我们用以接近这些主题的理智被愚蠢的观念和意志的乖戾所损伤和迟钝，让我们尽可能以谨慎的方式逐步获得对这件伟大之事的些许认识，其探究方式不像能看见的人，而像在黑暗中摸索的人。\par

% structure: p0006 source=h2 target=section reason=relative-heading-rank; base=h2

\section{第二章。——什么是恶。恶就是违反自然之物。承认这一点，摩尼派自相矛盾。}

2. 你们摩尼派常常——如果不是每次——问那些你们试图拉入你们异端的人：恶从哪里来？假设我现在是第一次遇见你们，我要请你们，如果你们愿意，暂时放下你们自以为对这些主题已经获得的解释，与我一起就像第一次一样开始这个伟大的探究。你们问我：恶从哪里来？我反过来问你们：什么是恶？哪个问题更合理？是那些不知道什么是恶就问它从哪里来的人正确，还是认为有必要先探究它是什么从而避免寻找未知之物起源的严重荒谬的人正确？当你们说恶就是违反自然之物时，你们的回答完全正确，因为没有人如此盲目以至于看不出，在每一种类中，恶就是违反该种类自然之物。但建立这一教义就是推翻你们的异端。因为如果恶违反自然，它就不是一个本质。然而，按照你们，恶是一个特定的本质和实体。此外，凡是违反自然的，必然反对自然并寻求毁灭它。因为\Quote{自然}无非意味着任何事物在其种类中被理解为存在的东西。因此才有了我们现在使用的新词，它来自表示存在的词——即\Quote{本质}或我们通常所说的\Quote{实体}——而在这些词使用之前，曾用\Quote{自然}一词。因此，如果你们不固执地思考这个问题，我们就会看到，恶就是那种脱离存在而趋向非存在的东西。\par

3. 因此，当天主教会宣告天主是一切本质和实体的创造者时，理解这一点的人同时也就理解天主不是恶的创造者。因为，那作为一切事物存在之因的，又怎能同时成为它们不存在——即脱离存在而趋向非存在——的原因呢？因为理性清楚地宣告这就是恶的定义。那么，你们那个被你们视为至恶的恶族，又怎能违反自然——即违反实体——如果按照你们，它本身就是一个本质和实体？因为如果它反对自身，它就毁灭自己的存在；当这完全实现时，它最终将成为至恶。但这是不可能的，因为你们不仅认为它存在，而且认为是永恒的。那么，那被说成是实体的东西，就不能是至恶。\par

4. 然而我该怎么做呢？我知道你们中许多人对此完全不能理解。我也知道有些人有很好的理解并能看见这些事，却如此固执地选择了恶——这种选择将也毁掉他们的理解——以致他们宁愿寻找能说些什么来蒙骗懈怠和软弱的心灵，而不愿认同真理。但我不会后悔写了这些，要么让你们中的某一天可能公正地思考它，并被引离开你们的错误，要么让那些理解且忠于天主、并且尚未被你们错误玷污的人读到它，从而得以免于被你们的言论引入歧途。\par

% structure: p0010 source=h2 target=section reason=relative-heading-rank; base=h2

\section{第三章。——如果恶被定义为有害之物，这暗示了对摩尼派的另一种驳斥。}

5. 让我们更仔细地探究，如果可能的话，更清晰地探究。我再问你们：什么是恶？如果你们说恶就是有害之物，那么在此你们也不会答错。但是，我请你们思考，我求你们警惕；请你们放下派别精神，为了寻找真理而探究，而不是为了战胜真理。无论什么有害之物，都会从它所危害的事物中剥夺某些善；因为没有善的损失，就不可能有伤害。我请问你们，还有什么比这更明白？还有什么更易于理解？对于一个中等理解力的人，只要他不乖僻，还需要什么才能完全证明这一点呢？但是，如果这一点被认可，那么后果似乎很明显。在你们认作首要之恶的那个族类中，没有任何东西可能受到伤害，因为它里面没有善。但是，如果正如你们所主张的，有两个本性——光明之国和黑暗之国；既然你们把光明之国当作天主，赋予它一个非复合的本性，以致它没有任何部分低于另一部分，那么你们必须承认——尽管这与你们自己强烈对立——你们必须承认，这个本性——你们不仅不否认它是首要之善，而且竭尽全力试图表明它是如此的——是不变的、不朽坏的、不可渗透的、不可侵犯的，否则它就不是首要之善；因为首要之善就是没有比它更好的善，而这样的本性是不可能受到伤害的。再者，如果，正如我们已经表明的，伤害就是剥夺善，那么对黑暗之国就不可能有伤害，因为它里面没有善。而且，既然光明之国不可伤害，因为它不可侵犯，那么你们所说的恶又对什么有害呢？\par

% structure: p0012 source=h2 target=section reason=relative-heading-rank; base=h2

\section{第四章 本身之善与分有之善的区别}

6. 现在，将你们无法逃脱的这种困惑与天主教会教导中的陈述的一致性相比较。根据这些教导，有一种善是至善的、本身是善，不是由于分有任何善，而是由于其本性和本质；另一种善是通过分有而成为善，由于接受了某种施予而成为善。因此，它作为善的存在来自于至善，而至善仍然是自足的，不失去任何东西。这第二种善被称为受造物，它因堕落而容易受到伤害。但天主不是这种堕落的创造者，因为祂是存在和存有的创造者。在这里，我们看到了恶这个词的正确用法；因为恶正确地应用于否定或损失，而非本质。我们也看到，什么样的本性容易受到伤害。这种本性不是首要之恶，因为当它受到伤害时，它失去善；它也不是首要之善，因为它从善堕落是因为它并非本身是善，而是由于拥有善。而一个事物不能在本性上是善，当它被说成是被造的时，这表明它的善是赋予的。因此，一方面，天主是善，祂所造的一切都是善，尽管不如造物主那样善。因为哪个疯子敢要求作品等同于工匠，受造物等同于创造者？你们还想要什么？你们还会希望有比这更明白的东西吗？\par

% structure: p0014 source=h2 target=section reason=relative-heading-rank; base=h2

\section{第五章 如果恶被定义为腐朽，这就彻底反驳了摩尼教异端}

7. 我第三次问：什么是恶？也许你们会回答：腐朽。不可否认，这是恶的一个普遍定义；因为腐朽意味着与本性的对立，也意味着伤害。但腐朽本身并不存在，而是存在于它所腐朽的某种实体中；因为腐朽本身并非实体。因此，它所腐朽的东西不是腐朽，不是恶；因为腐朽之物丧失了完整和纯洁。所以，没有纯洁可失去的东西不能被腐朽；而有纯洁的东西必然因分有纯洁而是善的。再者，腐朽之物是败坏的；败坏的东西丧失了秩序，而秩序是善的。因此，腐朽并不意味着善的缺失；因为在腐朽中它可以被剥夺善，如果善不存在，就不可能被剥夺。因此，那个黑暗之族，如果像你们所说的那样，缺乏一切善，就不可能被腐朽，因为它没有任何东西可以被腐朽夺走；如果腐朽不夺走任何东西，它就不腐朽。现在，如果你们敢，请说天主和天主的国可以被腐朽，因为你们无法表明魔鬼的国（如你们所设定的那样）如何能被腐朽。\par

% structure: p0016 source=h2 target=section reason=relative-heading-rank; base=h2

\section{第六章 腐朽影响什么与它是什么}

8. 公教之光还说什么？你猜是什么，但其实是真理：受造的实体是可败坏的，因为非受造的实体（即至善）是不败坏的；而败坏（即至恶）是不能败坏的，此外它也不是实体。但若你问败坏是什么，思考它试图将它所败坏的事物带向什么；因为它按照自己的本性影响那些事物。现在一切事物因败坏而脱离其本然，被引向非持续、非存在；因为存在意味着持续。因此，至上和至高的存在之所以如此称谓，是因为它持续于自身，或是自足的。在事物朝向更好的改变中，这种改变不是来自持续，而是来自对更坏者的背离，即从本质的脱离；这种脱离的作者不是本质的作者。所以在一些事物中有朝向更好的改变，从而趋向存在。这种改变不称为背离，而称为回归或转向；因为背离与有秩序的安排相对。现在，趋向存在的事物趋向秩序，获得秩序就获得存在，尽可能达到受造物所能达到的。因为秩序将所安排的事物化约为某种一致；存在无非就是成为一体。因此，某物获得多少一体性，它就存在多少。因为一致与和谐是一体性的效果，复合事物通过这些而存在，就它们所拥有的存在而言。因为单纯的事物自身存在，因为它们是一。而非单纯的事物通过其各部分的和睦模仿一体性，并且它们达到这一点，就存在这一点。这种安排是存在的原因，混乱是非存在的原因；背离或败坏是混乱的别名。所以凡是被败坏的事物都趋向非存在。现在可以让你自己去反思败坏的效果，以便你发现什么是至恶；因为它是败坏所要成就的。\par

% structure: p0018 source=h2 target=section reason=relative-heading-rank; base=h2

\section{第七章——天主的良善阻止败坏使任何事物归于无有。——创造与形成的区别。}

9. 但是天主的良善不允许这个结局实现，而是如此安排所有偏离的事物，使它们存在于它们最适宜存在的地方，直到在它们运动的秩序中返回它们所偏离的源头。因此，当理性灵魂远离天主时，虽然他们拥有最大限度的自由意志，天主仍将他们置于受造界的较低等级，那是他们应得的位置。于是他们因天主的审判而受苦，按他们的功过被安排到适当位置。由此我们看出你们一直强烈攻击的那句话的卓越性：\BibleQuote{我造善事，也造恶事。}\BibleRef{依 45:7} 创造就是形成与安排。所以在一些抄本中这样写道：\Quote{我造善事，并形成恶事。} 制造用于指先前不存在的事物；但形成是安排已有某种存在的事物，以改善和扩大它。这就是天主所说的\Quote{我形成恶事}，指那些正在偏离、因而趋向非存在的事物——不是那些已经达到它们所趋向的结果的事物。因为曾有人说过：\Quote{在天主的眷顾中，没有任何事物被允许走到非存在的尽头。}\par

10. 这些事可以更充分、更详细地讨论，但对你们来说，为达到我们的目的已经足够了。我们只需向你们展示你们绝望地寻找的门，也让未受教者对此绝望。你们只能通过善意进入，天主的仁慈将平安赐予善意的人，正如福音中的歌所唱：\BibleQuote{天主在至高之处受光荣，在地上平安归于善意的人。}\BibleRef{路 2:14} 我说，这足以向你们表明，解决善恶宗教问题的方法在于：凡是存在的，就其存在而言，都来自天主；而就它们偏离存在而言，则不来自天主，却又总是由天主眷顾按照整个体系安排。如果你们现在还不明白，我不知道还能做什么，只能更详尽地讨论已经说过的事。因为只有虔敬和纯洁才能引导心智走向更伟大的事物。\par

% structure: p0021 source=h2 target=section reason=relative-heading-rank; base=h2

\section{第八章——恶不是实体，而是敌视实体的不和谐。}

11. 因为对于\Quote{什么是恶？}这个问题，你还能给出什么其他答案？无非是：它反本性，或它有害，或它是败坏，或类似的说法？但我已经表明，在这样回答时，你们使你们的事业遭受覆灭，除非你们以通常对孩童讲话的那种幼稚方式回答，说恶是火、毒药、野兽，等等。因为这一异端的一位领袖（我们曾非常亲近和频繁地听他的教导），在谈到一个主张恶不是实体的人时，常说：\Quote{我想把一只蝎子放在那人手里，看他是否不会缩手；而他这样做，就得到了一种证明，不是用言语，而是用事实本身，证明恶是一种实体，因为他不会否认那动物是一种实体。} 他这样说并非当着那人的面，而是对我们说的，当时我们向他重复了那个令我们困扰的说法，如我所说，用幼稚的回答来回答幼稚的问题。因为稍有学识或科学素养的人，谁能不明白这些东西因与身体体质不合而伤害，而在其他时候又与体质相合，不仅不伤害，反而产生最好的效果？如果这毒药本身是恶的，那么蝎子自己首先会受最严重的伤害。事实上，如果毒药完全从动物身上移除，它就会死亡。因此，对它自己的身体而言，失去我们身体接受时有害的东西是恶的；而拥有我们缺乏时有益的东西是善的。那么，同一事物既是善又是恶吗？绝不是。恶是反本性的，因为这对于动物和对于我们都是恶。这种恶就是不合，它当然不是实体，而是敌视实体。那么它从哪里来？看它导致什么，你们就会知道，如果你们内里有任何内在之光的话。它把所有它毁灭的东西引向非存在。天主是存在的创造者；没有任何存在者，就其存在而言，引向非存在：因此我们知道了不合不是从哪里来的；至于它从哪里来，什么也不能说。\par

12. 我们在历史中读到，雅典有一名女囚，她分次服用了分配给死刑犯的致命剂量的毒药，却几乎或根本没有损害健康，因为她逐渐适应了毒药。因此，被判刑后，她像其他人一样按规定的剂量服毒，但因习惯化而使其无效，没有像其他人那样死去。此事引起了极大惊奇，于是她被放逐。如果毒药是恶，难道我们应该认为她使毒药对她不再是恶吗？还有什么比这更荒谬的呢？但因为不合是恶，她所做的只是通过习惯化的过程使有毒物质与她自己身体相合。因为无论多么狡猾，她怎能做到使不合不伤害她呢？为什么会这样？因为真正和恰当地称为恶的东西总是并且对所有人都有害。油对我们身体有益，但对许多六足动物却非常有害。难道藜芦有时是食物，有时是药物，有时是毒药吗？难道不是人人都承认，过量摄入盐是有毒的吗？然而盐对身体的益处不可胜数且极为重要。海水对陆上动物饮用有害，但对许多在其中沐浴身体者非常适宜和有用，对鱼类则两方面都有益且健康。面包滋养人，却杀死鹰。难道泥本身，在吞食或闻嗅时令人不快和有害，在热天却可用于冷却触摸，并治疗火伤吗？有什么比粪便更肮脏，比灰烬更无价值？然而它们对田地如此有用，以至于罗马人认为应对其发明者\Person{Stercutio}给予神圣的荣耀，\Quote{粪便}（\OriginalTerm{粪便}{stercus}）一词即源于其名。\par

13. 但为何列举数不胜数的细节？我们不必超出四大元素本身，众所周知，它们在被作用的对象中相合时有益，而相违时则与本性格格不入。我们生活在空气中，却死于土中或水中，而无数的动物在沙土或松土中活生生地爬行，鱼则死于我们的空气中。火消耗我们的身体，但在适当应用时，它既能从寒冷中恢复，又能驱除无数疾病。你们向之屈膝的太阳（诚然在可见事物中没有比它更美的），增强鹰的眼睛，但当我们注视它时，却伤害并模糊我们的眼睛；然而我们也能通过习惯化而毫无伤害地注视它。那么，你们愿意把太阳比作那位雅典女人通过习惯化使之无害的毒药吗？请反思一下，并考虑：如果一种实体因为伤害某人而是恶，那么你们所崇拜的光也不能免除这一指控。请看，把恶普遍地定义为这种不合是多么可取，太阳光线因其而产生眼睛的模糊，尽管没有比光更令眼睛愉悦的了。\par

% structure: p0025 source=h2 target=section reason=relative-heading-rank; base=h2

\section{第九章——摩尼教徒关于善恶的虚构自相矛盾。}

14. 我说这些话，是为了使你们——如果可能的话——不再称以下事物为\OriginalTerm{恶}{evil}：一个无垠深广的区域；一个游荡于区域中的\OriginalTerm{心灵}{mind}；五个\OriginalTerm{元素洞穴}{the five caverns of the elements}——一个充满\OriginalTerm{黑暗}{darkness}，另一个充满\OriginalTerm{水}{waters}，另一个充满\OriginalTerm{风}{winds}，另一个充满\OriginalTerm{火}{fire}，另一个充满\OriginalTerm{烟}{smoke}；以及生于各元素中的动物——黑暗中的\OriginalTerm{蛇}{serpents}，水中的\OriginalTerm{水生动物}{swimming creatures}，风中的\OriginalTerm{飞行动物}{flying creatures}，火中的\OriginalTerm{四足动物}{quadrupeds}，烟中的\OriginalTerm{两足动物}{bipeds}。因为按你们的描述，这些事物不能称为恶；因为所有这类事物，就其存在而言，必然从至高的天主获得存在，因为就其存在而言，它们是善的。如果\OriginalTerm{痛苦}{pain}和\OriginalTerm{软弱}{weakness}是恶，那么你们所说的那些动物具有如此强大的体力，以至于按照你们教派的信仰，当世界由它们形成之后，它们流产的后代从天上坠到地上，却不能死亡。如果\OriginalTerm{失明}{blindness}是恶，它们能看见；如果\OriginalTerm{耳聋}{deafness}是恶，它们能听见；如果近乎或完全\OriginalTerm{哑}{dumb}是恶，它们的言语如此清晰易懂，以至于按你们的主张，它们竟听从在集会上发表的演说决定与天主交战。如果\OriginalTerm{不育}{sterility}是恶，它们多产子嗣；如果\OriginalTerm{流放}{exile}是恶，它们身在本乡，占据自己的领土；如果\OriginalTerm{奴役}{servitude}是恶，它们中有些是统治者；如果\OriginalTerm{死亡}{death}是恶，它们活着，而且那种生命如此，按你们的说法，即使在天主得胜之后，心灵也不可能死亡。\par

15. 你们能告诉我，为何在\OriginalTerm{首恶}{the chief evil}中竟可发现如此多的善物，即上述诸恶的反面？如果这些不是恶，那么任何一个\OriginalTerm{实体}{substance}，只要是实体，能是恶吗？如果软弱不是恶，那么一个\OriginalTerm{软弱的身体}{weak body}能是恶吗？如果失明不是恶，那么黑暗能是恶吗？如果耳聋不是恶，那么\OriginalTerm{聋子}{deaf man}能是恶吗？如果哑不是恶，那么一条\OriginalTerm{鱼}{fish}能是恶吗？如果不育不是恶，我们怎能称一只\OriginalTerm{不育的动物}{barren animal}为恶？如果流放不是恶，我们怎能称一只\OriginalTerm{被流放的动物}{animal in exile}或使他人流放的动物为恶？如果奴役不是恶，那么一只\OriginalTerm{被征服的动物}{subject animal}或施行征服的动物在何种意义上是恶？如果死亡不是恶，那么一只\OriginalTerm{会死的动物}{mortal animal}或致死的动物在何种意义上是恶？或者，如果这些是恶，那么难道我们不应将体力、视力、听力、有说服力的言语、生育力、故乡、自由、生命这些你们都认为存在于那\OriginalTerm{恶的王国}{that kingdom of evil}中的事物称为善物，却竟敢称它为\OriginalTerm{恶的极致}{the perfection of evil}吗？\par

16. 再者，如果\OriginalTerm{不相宜}{unsuitableness}（从未被否认）是恶，那么还有什么比那些元素与其各自的动物更为相宜呢？——黑暗与蛇相宜，水与水生动物相宜，风与飞行动物相宜，火与贪婪的动物相宜，烟与高飞的动物相宜？这就是你们所描述的\OriginalTerm{争斗的族类}{the race of strife}中的\OriginalTerm{和谐}{harmony}；这就是\OriginalTerm{混乱之座中的秩序}{the order in the seat of confusion}。如果\OriginalTerm{有害的事物}{what is hurtful}是恶，我不重复已经提出的强烈异议——即若没有善存在，就不可能有伤害；但如果这一点不那么清楚，至少有一件事容易看出并理解，它来自公认的真理：有害的事物是恶。那区域中的烟没有伤害两足动物：它产生它们，滋养并维持它们，使它们的出生、成长和统治不受损伤。但现在，当恶与一些善混合后，烟变得更有害，以至于我们这些确实是两足动物的，非但不能靠它维持生命，反而被它弄瞎、窒息、杀死。难道善的混合会给恶元素带来如此大的毁灭性吗？难道天主的治理中会有如此混乱吗？\par

17. 至少在其他例子中，我们如何发现那种误导你们的作者、使他编造谎言的一致性？为什么黑暗与蛇一致，水与水生动物一致，风与飞行动物一致，尽管火会烧死四足动物，烟会呛死我们？再者，蛇不是视力极佳，喜爱阳光，并且在空气平静、云层很少聚集的地方最繁盛吗？多么荒谬啊，本是黑暗的土著和爱好者，竟然在最明亮的光线中生活得最舒适、最惬意！或者，如果你们说它们享受的是热而非光，那么将蛇（天生行动迅速）归给火，比归给行动迟缓的蝰蛇更合理。此外，所有人都必须承认，光对蝰蛇的眼睛是适宜的，因为它们被比作鹰的眼睛。但关于低级动物的话已经够了。我恳求我们注意自身真实的情况，不要坚持错误，这样我们的心灵就能摆脱愚蠢而有害的谎言。因为，如果说在\OriginalTerm{黑暗的族类}{the race of darkness}中，在没有光混合的情况下，\OriginalTerm{两足动物}{biped animals}具有如此健全、强壮、令人难以置信的视力，以至于即使在他们的黑暗中也能看见（如你们所描述的）\OriginalTerm{天主的国}{the kingdom of God}的完全纯净的\OriginalTerm{光}{light}，这岂非难以容忍的悖逆吗？因为按照你们的说法，就连这些存在也能看见这光，凝视它，研究它，喜爱它，渴求它；而我们的眼睛，在与光、与\OriginalTerm{至善}{the chief good}、甚至与天主混合之后，却变得如此柔嫩脆弱，以至于我们既不能在黑暗中看见任何东西，也不能直视太阳，而且看后还会失去之前能看见的视力。\par

18. 若我们把腐败视为一种恶（无人怀疑），同样的评语亦适用。烟雾并未败坏那群动物，尽管它如今败坏动物。不必逐一列举所有细节——那将冗长，且无必要——，你虚构描述中的生物远比现今的动物不易腐败，以致它们那从天上摔到地上的早产、未成熟的后代，竟能存活、繁殖，并能重新聚合，因它们有（据称）原有的活力，因为它们在善与恶混合之前就已受孕；因为照你们说，经过这混合之后，所生的动物便是我们如今所见的那些十分软弱、易于腐败的生物。有谁能在这种错误信念中坚持下去呢？除非他看不见这些事，或受你们的习惯和交往影响得如此惊人，以致对一切推理的力量都无动于衷？\par

% structure: p0031 source=h2 target=section reason=relative-heading-rank; base=h2

\section{第十章——摩尼教徒无端订立的三项道德标记}

19. 如今，我已指出（如我所想）你们的意见中关于一般善与恶的事物有多少黑暗与错误，让我们现在来审视你们大肆赞扬并自诩为卓越实践的那三项标记。那么，这三项标记是什么呢？口的标记、手的标记、胸的标记。这是什么意思？我们被告知，人应在口、手、胸上保持纯洁无辜。但如果他用眼、耳、鼻犯罪呢？如果他用了脚跟伤害某人，甚或杀了他呢？他怎能被算为有罪，当他没有用口、手、胸犯罪时？但有人回答说：口，我们理解为头部所有感官器官；手，一切身体行为；胸，一切情欲倾向。那么，你将亵渎归于何处？归于口还是手？因为亵渎是舌头的行动。如果一切行为都归入同一类别，何以你们将手和脚的行为合在一起，而不将舌头的行为也包括进去？你们是否想将舌头的行动（因其意在表达某事）与不意在表达的行动分开，以便手的标记意味着戒除一切不意在表达的恶行？但这样，如果有人用手表达某事而犯罪，如书写或某种有意义的手势，那该如何？这不能归于舌头和口，因为它是由手做的。当你们有口、手、胸三项标记时，全然不能接受将手上的罪归于口。而如果你们将一般行动归于手，那也没有理由只包括脚的行动而不包括舌头的行动。你们可看到，对新奇的渴求及其伴随的错误，将你们引入多大的困境？因为你们发现不可能将一切罪的净化包含在这三项标记中，而你们却将它们提出作为一种新的分类。\par

% structure: p0033 source=h2 target=section reason=relative-heading-rank; base=h2

\section{第十一章——摩尼教徒中口的标记的价值，他们被发现犯有亵渎天主之罪}

20. 随意分类，随意省略吧；我们必须讨论你们最坚持的教义。你们说，口的标记意味着戒除一切亵渎。但亵渎就是说好东西的坏话。所以\Quote{亵渎}一词通常只用于说天主的坏话；因为关于人，尚有不确定性，但天主毫无争议是善的。那么，如果你们被证明说了比别人更坏的天主的话，你们那著名的口的标记又成什么了？证据并不隐晦，而是清晰明显、无人能抗拒；且无论何人对此都不能无知，即天主是不朽的、不变的、不受任何伤害、没有任何匮乏、没有软弱、没有痛苦。这一切理性受造物的一般知觉都能觉察，甚至你们听到时也同意。\par

21. 但当你们开始讲述你们的寓言，说天主是可朽的、可变的、会受伤害、会匮乏、软弱、且难免痛苦时，这就是你们盲目地教导、有些人盲目地相信的东西。这还不算；按你们所说，天主不仅可朽，而且已被败坏；不仅可变，而且已被改变；不仅会受伤害，而且已被伤害；不仅可能匮乏，而且已经匮乏；不仅可能软弱，而且确实软弱；不仅会痛苦，而且已经痛苦。你们声称灵魂就是天主，或天主的一部分。我不明白，部分的天主如何能不是天主？金的一部分是金；银的一部分是银；石的一部分是石；并且（来到更大的事物上）土的一部分是土，水的一部分是水，气的一部分是气；如果你从火中取一部分，你不会否认那是火；光的一部分只能是光。那么，为什么天主的一部分就不会是天主？难道天主有像人及低级动物一样的关节身体？因为人的一部分并不是人。\par

22. 我将分别处理这些意见。你若视天主为光的模样，就必须承认天主的一部分即是天主。因此，当你使灵魂成为天主的一部分时，尽管你允许它因愚妄而败坏、因曾明智而改变、因需要健康而匮乏、因需要医治而衰弱、因渴望幸福而悲惨，你便亵渎地将这一切归于天主。或者，你若否认灵魂有这些属性，那么圣神就不必引导灵魂进入真理，因为它不处于愚妄之中；灵魂也不必因真宗教而更新，因为它不需要更新；它也不必藉你们的象征而达于完美，因为它已然完美；天主也不必给予它帮助，因为它不需要；基督也不必作它的医师，因为它健康；它也不必在另一生命中得享幸福的应许。那么，为何耶稣被称为解救者，正如祂在福音中所说：\BibleQuote{若子使你们自由，你们就真正自由了。}（\BibleRef{若 8:36}）而宗徒保禄也说：\BibleQuote{你们蒙召是为了自由。}（\BibleRef{迦 5:13}）因此，尚未获得这自由的灵魂仍处于束缚之中。那么，按照你们的说法，天主（因为天主的一部分即是天主）既因愚妄而败坏，又因堕落而改变，因失去完美而受损，因需要帮助而匮乏，因疾病而衰弱，因悲惨而屈辱，且受可耻的束缚。\par

23. 再者，若天主的一部分并非天主，那么当祂的部分被败坏时，祂仍非不朽；当祂的任何部分改变时，祂仍非不变；当祂在每一部分并非完美时，祂仍非无损；当祂忙于恢复自身的一部分时，祂仍非无缺；当祂有软弱的部分时，祂仍非全然健全；当任何部分遭受悲惨时，祂仍非完全幸福；当任何部分处于束缚时，祂仍非完全自由。这些是你们被迫得出的结论，因为你们说灵魂（你们看见它处于如此灾难性的状态）是天主的一部分。若你们能使你们的教派放弃这些及许多类似的意见，那么你们才可以说自己的口远离了亵渎。更好的是，离开该教派；因为若你们不再相信并重述摩尼所写的一切，你们就不再是摩尼教徒了。\par

24. 天主是至善，是比它更善者无法被想象或存在的，我们必须如此理解或相信，如果我们希望避免亵渎。有一种数字关系，它不可能被损害或改变，任何本性无论以何种暴力都不能阻止\Quote{一之后}的数字成为\Quote{一}的两倍。这无论如何也无法改变；而你们却把天主描述为可变的！这种关系保持其完整性不受侵犯，而你们甚至不允许天主在这方面拥有平等！让某个黑暗种族抽象地取数字三（由不可分的单位组成），并将其等分为两部分。你的心智明白，没有任何敌意能实现这一点。那么，那不能损害数字关系的东西，怎能损害天主呢？若它不能，怎么可能有必要让天主的一部分与邪恶混合，并被驱入如此悲惨的境地呢？\par

% structure: p0039 source=h2 target=section reason=relative-heading-rank; base=h2

\section{第十二章——摩尼教的遁词。}

25. 由此产生了一个问题，即使当我们还是你们的热心门徒时，这个问题也常使我们陷入极大的困惑，而我们找不到任何答案——假若天主拒绝以自身一部分遭遇如此灾难的代价与黑暗种族作战，那黑暗种族会如何对待天主？因为若天主保持安静不会遭受任何损失，我们便认为我们被派去忍受如此多苦难是苛刻的。再者，若祂会遭受损失，祂的本性便不可能是不朽的，而天主的本性应当是不朽的。有时有人回答说：并非为了逃避邪恶或避免伤害，而是天主出于祂固有的良善，愿意将秩序赐予一个混乱无序的本性。但我们在摩尼教的书中并未发现这一点；那些书始终暗示并断言天主防备了敌人的入侵。但假定这一答案（因无更好的）代表了摩尼教的意见，那么天主是否就免于残酷或软弱的指控了呢？因为祂对敌对种族的这种良善对祂自己的子民却极为有害。此外，若天主的本性既不能被败坏也不能被改变，那么任何毁灭性的影响也不能败坏或改变我们；而赐予陌生种族的秩序，本可以不剥夺我们的秩序而赐予。\par

26. 然而，自那时起，另一个答案出现了，我最近在迦太基听到的。因为有一位我极希望看到他脱离此错误的人，当他被逼到同一困境时，竟敢说王国有自己的界限，可能被敌对种族侵犯，尽管天主自身不能受伤害。但你们的创始人绝不会同意这个回答；因为他很可能看出这种意见会更快地摧毁他的异端。事实上，任何中等智力的人，听到在这种本性中一部分易受伤害、一部分不易受伤害，会立刻看出这构成了不是两种而是三种本性：一种可侵犯的，一种不可侵犯的，第三种侵犯的。\par

% structure: p0042 source=h2 target=section reason=relative-heading-rank; base=h2

\section{第十三章——行为的判断应基于动机而非外表。以这一原则检验摩尼教的斋戒。}

27. 你们口中天天有这些发自内心的亵渎，你们不该继续高举口的象征作为奇妙之物，去欺骗无知的人。但也许你们认为口的象征卓越而可敬，因为你们不吃肉、不喝酒。但你们的目的是什么？因为根据我们行动时所着眼的目的（即我们做一切事的原因），不仅不是可指责的，反而是可称赞的，只有如此我们的行动才能得任何赞美。若我们在任何行为中所考虑的目的是不合法、该受责备的，那么该行为本身就会被毫不犹豫地谴责为不当。\par

28. 我们听说\Person{喀提林}能忍受寒冷、口渴和饥饿。这个卑鄙的恶棍在这点上与我们的宗徒相同。那么，是什么将弑亲者与我们的宗徒区分开来，难道不是他所追求的截然相反的目的吗？他忍受这些是为了满足他凶猛而无法控制的激情；而宗徒们则是为了克制这些激情，使之服从理性。当你们听说有许多天主教贞女时，你们常说：\Quote{一头母骡也是贞女。}这确实是对天主教制度的无知说法，并不适用。尽管如此，你们的意思是，这种节制除非以正确的原则导向卓越的终极，否则毫无价值。天主教基督徒也可以将你们戒酒戒肉的行为与牛、许多小鸟以及无数种虫子相比。但是，为了不像你们那样无礼，我不会过早地进行这种比较，而是首先审视你们行为的目的。因为我想我可以稳妥地认为大家同意：在这样的习俗中，目的是要关注的。因此，如果你们的目的在于节俭和抑制在饮食中寻求满足的欲望，我赞同并认可。但事实并非如此。\par

29. 假设可能存在这样一个人：他在饮食上非常节俭，非但不满足食欲或口腹之欲，反而节制自己一天只吃两餐，晚餐只吃少许用猪油润泽调味的卷心菜，刚刚够抑制饥饿；为了健康，他喝两三杯纯酒解渴；这就是他日常的饮食。而另一个人习惯不同，从不吃肉或酒，却在下午两点享用一顿由稀有的异国蔬菜组成的愉快餐食，菜品多样，撒满胡椒；到了晚上也以同样方式进餐；他喝蜜醋、蜂蜜酒、葡萄干酒以及各种水果的汁液，这些饮品很好地模仿了酒，甚至在甜度上超过它；他喝酒不是为了解渴，而是为了享乐；他每天为自己准备这些，如此丰盛地吃喝，不是出于需要，而仅仅是为了满足味觉。在这两人中，你认为谁在饮食上最为节制？你绝不可能盲目到不把那个吃少许猪油和喝点酒的人放在这个饕餮之徒之上！\par

30. 这才是正确的看法；而你们的教义听起来却大相径庭。因为你们的一位以三个标记著称的选民可以像上述描述中的第二个人那样生活，尽管他可能会被一两个较为持重的人责备，但他不能因滥用标记而被定罪。但是，如果他和另一个人一起吃饭，用一块腐臭的咸肉润湿嘴唇，或者用一杯变质的酒解渴，他就被宣布为违反了标记，根据你们创始人的判定被投入地狱，而你们虽然惊讶，却必须赞同。你们难道不抛弃这些错误吗？你们难道不听理性的话吗？你们难道不对习惯的力量作一点抵抗吗？这样的教义难道不是极其不合理吗？这难道不是疯狂吗？一个人每天用蘑菇、大米、块菌、蛋糕、蜂蜜酒、胡椒和阿魏胶填满肚子以满足食欲，并且每天都这样吃喝，却不能因违反三个标记（即圣洁的规则）而被定罪；而另一个人用一些烟熏的肉块调味最普通的蔬菜，只摄取身体恢复所需的分量，为了保持健康只喝三杯酒，却因将前者的饮食换成这个而被判受罚，这难道不是最大的荒谬吗？\par

% structure: p0047 source=h2 target=section reason=relative-heading-rank; base=h2

\section{第十四章——戒除某些食物的三个好理由}

31. 但是，你们回答说，宗徒说：\BibleQuote{弟兄们，不食肉、不饮酒，这是好的。}（\BibleRef{罗 14:21}）没有人否认这是好的，只要是为了已经提到的目的，即如经上所说：\BibleQuote{不要为肉体安排，去放纵私欲。}（\BibleRef{罗 13:14}）或者为了宗徒所指出的目的：或是抑制食欲（食欲在这些东西上容易比在其他东西上更加无节制地放纵），或是免得弟兄跌倒，或是免得软弱的人与偶像有分。因为宗徒写作的时候，祭肉常常在市场上出售。而且，由于酒也用于外邦人的奠祭，许多软弱的弟兄习惯于购买这些东西，宁愿完全戒除肉和酒，也不愿冒险——即使是无知地——与偶像有分，他们认为那是与偶像相交。为了他们的缘故，甚至那些更强壮、有足够信心看透这些事微不足道的人，知道没有什么是不洁的，除非出于邪恶的良心，并持守主的这句话：\BibleQuote{入口的不能污秽人，出口的才能污秽人。}（\BibleRef{玛 15:2}）但仍然必须戒除这些东西，免得这些软弱的弟兄跌倒。这不是单纯的理论，而是宗徒在书信中明确教导的。因为你们习惯于只引用\BibleQuote{弟兄们，不食肉、不饮酒，这是好的}，却不加上后面的话：\Quote{也不作任何使你弟兄跌倒、或绊倒、或软弱的事。}这些话显示了宗徒给予劝诫的意图。\par

32. 从前后文可以清楚看出这一点。这段经文很长，但为了那些不勤于阅读和查考圣经的人，我们必须完整地引用它。宗徒说：\Quote{信德软弱的，你们要接纳，但不要辩论所疑惑的事。有人信凡物都可吃；但那软弱的，只吃蔬菜。吃的人不可轻看不吃的人；不吃的人不可论断吃的人，因为天主已经收纳他了。你是谁，竟论断别人的仆人？他或站住或跌倒，自有他的主人在；而且他也必站住，因为主能使他站住。有人看这日比那日强；有人看日日都一样。只是各人心里要意见坚定。守日的人，是为主守的；吃的人，是为主吃的，因他感谢天主；不吃的人，是为主不吃的，也感谢天主。我们没有一个人为自己活，也没有一个人为自己死。我们若活着，是为主而活；若死了，是为主而死；所以我们或活或死，总是主的人。为此基督死了，又活了，为要作死人并活人的主。你为甚么论断弟兄呢？又为甚么轻看弟兄呢？因我们都要站在天主的审判台前。经上记着说：主说：我指着我的永生起誓，万膝必向我跪拜，万口必向天主承认。\BibleRef{依 45:23}这样看来，我们各人必要将自己的事在天主面前说明。所以我们不可再彼此论断；宁可定意谁也不给弟兄放下绊脚跌人之物。我凭着主耶稣确知深信，凡物本来没有不洁净的；惟独人以为不洁净的，在他就不洁净。你若因食物叫弟兄忧愁，就不是按着爱德而行。基督已经替他死，你不可因你的食物叫他败坏。不可叫你的善被人毁谤。因为天主的国不在乎吃喝，只在乎公义、和平，并圣神中的喜乐。凡这样事奉基督的，就为天主所喜悦，又为人所称许。所以我们务要追求和睦的事与彼此建立德行的事。不可因食物毁坏天主的工程。凡物固然洁净，但有人因食物叫人跌倒，就是他的罪了。无论是吃肉，是喝酒，是甚么别的事，叫弟兄跌倒，一概不作才好。你有信德，就当在天主面前守着。人在自己以为可行的事上能不自责，就有福了。若有疑心而吃的，就必有罪，因为他吃不是出于信德；凡不出于信德的都是罪。我们坚固的人应该担代不坚固人的软弱，不求自己的喜悦。我们各人务要叫邻舍喜悦，使他得益处，建立德行。因为基督也不求自己的喜悦。}\par

33. 难道不清楚吗？宗徒所要求的是，强壮的不应吃肉喝酒，因为他们不同意软弱的人，就绊倒了软弱的人，并且使那些凭信德判断一切洁净的人，以为不回避这类食物饮料是在敬拜偶像？这一点在\WorkTitle{格林多前书}的以下段落中也显明了：\Quote{论到吃祭偶像之物，我们知道偶像在世上算不得甚么，也知道天主只有一位，再没有别的神。虽有称为神的，或在天上，或在地下，就如那许多的神，许多的主；然而我们只有一位天主，就是圣父，万物都本于他，我们也归于他；并有一位主，就是耶稣基督，万物都是藉着他有的，我们也是藉着他有的。但人不都有这等知识；有人到如今因拜惯了偶像，就以为所吃的是祭偶像之物；他们的良心既然软弱，也就污秽了。其实食物不能叫天主看中我们，因为我们不吃也无损，吃也无益。只是你们要谨慎，恐怕你们这自由竟成了那软弱人的绊脚石。若有人见你这有知识的在偶像的庙里坐席，这人的良心，若是软弱，岂不放胆去吃那祭偶像之物吗？因此，基督为他死的那软弱弟兄，也就因你的知识沉沦了。你们这样得罪弟兄们，伤了他们软弱的良心，就是得罪基督。所以，食物若叫我弟兄跌倒，我就永远不吃肉，免得叫我弟兄跌倒。}\BibleRef{格前 8:4}\par

34. 又在另一处：\Quote{我是怎么说呢？是说祭偶像之物算得甚么呢？或说偶像算得甚么呢？我乃是说，外邦人所献的祭是祭鬼，不是祭神；我不愿意你们与鬼相交。你们不能喝主的杯，又喝鬼的杯；不能吃主的筵席，又吃鬼的筵席。我们可惹主的愤恨吗？我们比他还有能力吗？凡事都可行，但不都有益处；凡事都可行，但不都造就人。无论何人，不要求自己的益处，乃要求别人的益处。凡市上所卖的，你们只管吃，不要为良心的缘故问甚么话；若有人对你们说：这是献过祭的物，就要为那告诉你们的人，并为良心的缘故不吃。我说的良心，不是你的，乃是他的。我这自由为甚么被别人的良心论断呢？我若谢恩而吃，为甚么因我谢恩的物被人毁谤呢？所以，你们或吃或喝，无论做甚么，都要为荣耀天主而行。不拘是犹太人，是希腊人，是天主的教会，你们都不要使他跌倒；就好像我凡事都叫众人喜欢，不求自己的益处，只求众人的益处，叫他们得救。你们该效法我，像我效法基督一样。}\par

35. 那么，我认为很清楚，我们应当为了什么目的而戒绝肉食和酒。这目的有三：为了节制放纵，这类食物大多使用放纵，这类饮品则会导致酩酊大醉；为了保护软弱，因为那些祭献和奠祭之物；以及最值得称赞的，出于爱德，不使比我们更软弱的、对这些事物持戒的人跌倒。而你们，却认为一口肉是不洁的；但宗徒却说万物都是洁净的，但那带着疑虑吃的人就有罪了。无疑地，你们被这样的食物玷污，仅仅是因为你们认为它不洁。因为宗徒说：\BibleQuote{我凭着主耶稣确知深信，凡物本来没有不洁净的；惟独人以为不洁的，在他就不洁了。}（\BibleRef{罗 14:14}）每个人都能看出，他所说的\Quote{不洁}即是指不洁和污秽。但和你们讨论圣经段落是愚蠢的；因为你们一方面许诺要证明你们的教义来误导人，另一方面又断言那些有权威要求我们崇敬的书籍被篡改的赝品所败坏。那么，请向我证明你们的教义：肉食玷污食者，当它是在不冒犯任何人、没有软弱观念、没有过度的情形下被食用时。\par

% structure: p0053 source=h2 target=section reason=relative-heading-rank; base=h2

\section{第十五章——摩尼教徒为何禁止食用肉食}

36. 值得注意他们迷信禁食的全部理由，如下所述：——我们被告知，既然天主的肢体已经与恶之实体混合，以克制它并防止它过分凶暴——这就是你们所说的——世界是由两种本性（善与恶）混合而成的。但这部分天主每日在世界各处被释放，并归还到其本域。但它以蒸汽形式从大地上升至天堂时，进入植物，因为植物的根扎在土中，从而赋予所有草木和灌木以肥力和力量。动物从这些植物获得食物，而在有性交的地方，它们将天主的肢体禁锢在肉体中，使其偏离正道，并妨碍它、缠绕它于错误和困苦中。因此，如果由蔬菜和水果组成的食物来到圣人（即摩尼教徒）那里，通过他们的贞洁、祈祷和圣咏，其中任何卓越和神圣的部分都被净化，从而完全成全，以便毫无阻碍地归还到其本域。因此，你们禁止人们把面包、蔬菜甚至水（这些东西不费任何代价）施舍给乞丐，如果他不属于摩尼教徒，以免他因自己的罪玷污了天主的肢体，并阻碍其回归。\par

37. 你们说，肉食本身就是由污秽构成的。因为按照你们，那神圣部分有一些在食用蔬菜和水果时逃脱了：它逃脱于它们遭受擦磨、碾碎或烹煮的过程，也逃脱于咬嚼的动作。它也逃脱于动物的所有运动中，在负重、操练、劳苦或任何行为中。它也逃脱于我们的休息中，当消化通过体内热量进行时。由于神圣本性以所有这些方式逃脱，就留下了一些极其不洁的渣滓，从这些渣滓中，通过性交形成肉食。然而，这些渣滓在上述运动中，连同灵魂中的善一起飞散；因为它虽大多是善，但不完全是善。所以，当灵魂离开肉食后，渣滓就极其污秽，而食用肉食者的灵魂就被玷污了。\par

% structure: p0056 source=h2 target=section reason=relative-heading-rank; base=h2

\section{第十六章——揭露摩尼教徒的怪诞教义}

38. 哦，事物本性的晦暗！揭露虚假是多么困难！凡听见这些话的人，如果他没有学习过事物的原因，尚未被任何真理之光启迪，却被物质形象所欺骗，难道不会认为它们是真的吗？正是因为所说的事是看不见的，并以可见事物的形式呈现于心灵，而且能被雄辩地表达出来。这样的人数量众多，成群结队；他们未被这些错误引诱，更多是出于宗教情感所激发的敬畏，而非理性的推理。因此，我将尽力，按天主所愿赐予我的能力，来驳斥这些错误，使得它们的虚假和荒谬不仅显明于智者的判断——他们一听就拒绝——也显明于常人的理解力。\par

39. 那么，首先告诉我，你们从哪里得到这个教义：天主的某一部分（如你们所说）存在于谷物、豆子、卷心菜、花朵和果实中？他们说，从颜色的美丽和滋味的甘甜可以显明这一点；由于腐烂的物质中没有这些，我们知道它们的善已被剥夺。他们岂不羞愧将发现天主归功于鼻子和味觉吗？但我暂且不说这个。因为我要使用词语的本义来讲；正如俗语所说，对你们讲话并非那么容易。我们不如看看，需要什么样的心智才能理解这一点：如果善在身体中的临在由其颜色显示，那么动物的粪便、肉食本身的残渣都有各种明亮的颜色，有时白色，常常金色，等等；而这些正是你们在果实和花朵中作为天主临在和内住的证据。为什么在玫瑰中，你们将红色视为善丰裕的标记，而在血液中同样的红色却遭到你们谴责？为什么你们在紫罗兰中愉快地看待同一种颜色，却在霍乱、黄疸病人或婴儿的排泄物中避之不及？为什么你们相信油的闪亮外观是善大量混合的标志，你们乐于通过将油吸入喉和胃来净化它，却不敢用脂肪的一滴触碰嘴唇，尽管它有与油一样的闪亮外观？为什么你们将黄色的甜瓜视为天主宝藏的一部分，而将腐臭的培根脂肪或蛋黄排除在外？为什么你们认为生菜中的白色宣告天主，而牛奶中的白色却不？关于颜色就谈这么多（且不提别的），你们无法将任何鲜花盛开的草地与一只孔雀的翅膀和羽毛相比，尽管这些都属于肉食且源于肉食。\par

40. 再者，如果这善也能通过嗅觉被发现，那么用某些动物的肉制成的高级香水就具有极佳的香味。而且，食物与味道鲜美的肉一起烹调时，其气味比不加肉时更好。再者，如果你认为气味比别的东西更好的东西因此就更洁净，那么有一种泥巴你应该带进餐中，而不是从水窖中取水；因为被雨水湿润的干土有一种最令人愉悦的气味，而且这种泥巴比单独的雨水气味更好。但如果我们必须依靠味觉的权威来证明任何物体中有一部分天主，那么祂必更多地存在于枣子和蜂蜜中，而非猪肉中，但更多存在于猪肉中而非豆类中。我承认祂更多地存在于无花果中而非肝脏中；但这样你就必须承认祂更多地存在于肝脏中而非甜菜中。而且，按照这个原则，你岂不必须承认有些植物——你们无人能否认它们比肉更洁净——是从这肉本身接收天主的，如果我们认为天主与味道相混合的话？因为卷心菜与肉一起烹调时味道更好；而且，虽然我们无法喜欢牲畜吃的植物，但当这些植物变成奶时，我们认为它们的颜色改善了，并且觉得它们非常可口。\par

41. 还是我们必须认为，善存在于三种善的品质——好的颜色、气味和味道——同时具备的地方更多？那么你们就不应过分赞美花，因为你们不能接受让它们接受味觉的审判。至少你们不应把马齿苋看得比肉更高，因为肉在烹调后颜色、气味和味道都更优。一只烤乳猪（因为你们在这问题上的观点迫使我们像对待厨师和糕点师一样与你们讨论善恶，而不是读书人或文学品味的人）颜色明亮，气味宜人，味道可口。这是神圣实体存在的完美证据。你们被这三重见证邀请，并被召唤通过你们的圣洁来净化这一实体。动手吧。你们为何退缩？你们有什么反对意见？仅就颜色而言，婴儿的排泄物胜过扁豆；仅就气味而言，一块烤肉胜过柔软的青无花果；仅就味道而言，一只被杀的小山羊胜过它活着时吃的植物：我们已经发现一种肉，其味道这三种证据都具备。你们还要求什么？你们会如何回答？如果食用这样的怪物进行讨论不会使你们不洁，那么为什么吃肉会使你们不洁？最重要的是，太阳的光线——你们肯定认为它比所有动植物食物更高——没有气味或味道，仅以极其明亮的颜色在其他物质中突出；这大声呼唤你们，并迫使你们，尽管你们不愿意，也绝不可将任何东西置于亮色之上作为善之混合的证据。\par

42. 因此你们被迫陷入这样的困境：你们必须承认天主更多地居住在血中，以及丢弃在街上的肮脏但颜色明亮的动物粪便中，而非橄榄的苍白叶子中。如果你们回答——正如你们实际所做的——橄榄叶燃烧时发出火焰，证明光的存在，而肉燃烧时却不如此，那么对于油——它几乎点亮了\Place{意大利}所有的灯——你们会说什么？对于牛粪（它肯定比肉更不洁净），农民们晒干后用作燃料，以致火总是即时可得，烟尘不断释放，你们又说什么？如果明亮和光泽证明神圣部分的存在更多，那么你们自己为何不净化它？为何不占有它？为何不解放它？因为它主要存在于花中，更不用说血以及肉中或来自肉的无数几乎与血相同的东西了，而你们却不能以花为食。即使你们吃肉，你们也肯定不会把鱼鳞、或者一些蠕虫和苍蝇与你们的稀粥一起吃，尽管这些东西在黑暗中都有它们自己的光泽。\par

43. 那么，还剩下什么？只有停止说你们在眼睛、鼻子和味觉中有足够的手段来测试神圣部分在物质对象中的存在。而没有这些手段，你们如何能辨别——不仅在植物中比在肉中有更多的天主部分，而且植物中究竟是否有任何部分？你们是否被它们的美丽——不是令人愉悦的颜色的美，而是各部分协调的美——引导而这样想？依我之见，这是一个极好的理由。因为你们永远不会如此大胆，以至于将扭曲的木块与由彼此相称的肢体构成的动物身体相比。但如果你们选择感官的见证——就像那些不能用心灵看到存在之全部力量的人必须做的那样——你们如何证明善的实体随着时间的推移并从某种摩擦中逃离物体，但因为在你们看来，天主已经离开了它，并离开了一个地方而去了另一个地方？整个说法是荒谬的。但是，就我所能判断的，没有任何标志或现象可以引起这种意见。因为许多从树上摘下或从土里拔出的东西，在食用前经过一段时间反而更好，例如韭葱、菊苣、生菜、葡萄、苹果、无花果和一些梨；还有许多其他东西，在摘下后不立即使用时反而颜色更好，而且对身体更有益，对味觉更美味。但这些东西不应该具有所有这些优良和令人愉悦的品质，如果如你们所说，它们离开母土后保存得越久，就越缺乏善。动物肉本身在动物被杀后的第二天味道更好、更适合食用；但这是不应该的，如果如你们所坚持的，它在屠宰后立即比第二天拥有更多的善，因为那时更多的神圣实体已经逃逸。\par

44. 谁不知道陈酒更纯净、更佳美？它并非如你们所想的那样更能诱惑人败坏感官，而是更有益于强身健体——只要节制，这应当统御一切。新酒反而更快地摧毁感官。新酒在酒桶中仅短暂时日，开始发酵时，便使俯视桶中的人头晕目眩，以致若不扶持便会丧命。至于健康，谁不知道新酒会使身体肿胀并有害膨胀？难道只因其中含有更多善质，它就有这些不良特性吗？这些特性在陈酒中不出现，岂非因为大量神圣实体已经离去？这种说法荒谬绝伦，尤其对你们而言，你们正是以眼、鼻、舌所感受到的快感来证明神圣临在的！你们一面把酒说成黑暗首领的毒药，一面又吃葡萄，这是何等矛盾！难道酒在杯中比在葡萄串中具有更多毒药？或者，若善离去后恶仍不混合，且随时间推移，为何同一葡萄悬挂一段时间后变得更温和、更甘甜、更健康？又为何如上所述，酒本身在光明离去后变得更纯净澄清，因失去有益实体而更健康？\par

45. 至于木材和叶子，它们随时间变干，但在你们的评判中不应因此变得更糟吧？因为它们既失去了产生烟雾之物，又保留了产生明亮火焰之物；且以你们所看重的清澈而论，干木比青木有更多善。因此，你们要么否认纯光中比烟光中有更多天主——这将推翻你们所有的证据；要么就得承认，植物被拔起或树枝被折下并放置一段时间后，从它们中逸出的恶的本性可能多于善的本性。据此，我们就会认为从摘下的水果中可能逸出更多恶；因此动物性食物中可能保留更多善。关于时间的问题就谈到这里。\par

46. 至于运动、抛掷和摩擦，若这些给予神圣本性逃离这些物质的机会，那么有许多同类事物对你们不利，它们因运动而改善。某些谷粒的汁液类似酒，经搅拌后极佳。事实上，不可忽视的是，这类饮料迅速致醉；然而你们从未称谷粒的汁液为黑暗首领的毒药。有一种用水加少许面粉调成的制剂，摇动后更好，奇怪的是，光明离去后颜色反而变浅。糕点师长时间搅动蜂蜜使其颜色变浅，甜味更柔和、少不健康：你们必须解释这如何来自善的损失。再者，若你们更喜欢以听觉而非视觉、嗅觉或味觉的快感来检验天主的临在，那么竖琴的弦来自动物的肠，风笛的管来自动物的骨；它们经干燥、摩擦和扭转才变得有音乐性。因此，你们认为来自神圣王国的音乐之乐，竟是从死动物的废弃物中获得的，且经时间干燥、摩擦减少、扭转拉长。按照你们的说法，这种粗暴处理会驱使神圣实体离开活物；你们甚至说烹饪也是如此。为何煮过的蓟菜反而无害？难道因为天主或天主的一部分在烹饪时离开了它们？\par

47. 何必一一列举？要一一列举也很困难，也没有必要；因为谁都知道有多少东西煮熟后更甘甜、更健康。若如你们所猜测，事物因这样运动而失去善，这就不应该发生。我不认为你们能凭身体感官证明肉是不洁的，会玷污食者灵魂，因为水果经采摘并以各种方式摇动后成为肉；尤其你们认为醋因陈化发酵比酒更纯净，而你们饮用的蜂蜜酒不过是煮熟了的酒，若物质因被运动、烹饪而失去神圣肢体，则蜂蜜酒应比酒更不洁。但若不然，你们就没有理由认为水果经采摘、存放、处理、烹饪、消化后，被善所遗弃，从而为身体的形成提供最不洁的材料。\par

48. 但若你们不是凭颜色、外观、气味和味道而认为善在这些事物中，那你们还能提出什么？你们是否从这些事物离开土地并被使用后似乎失去的力量和活力来证明？若这是你们的理由（尽管其错误显而易见，因有些事物离开土地后力量反而增强，如前所述，酒经陈年而更强）——那么若以力量为理由，则天主的部分在肉类中比在任何食物中都更丰富。因为尤其需要力量和精力的运动员，并不惯于以卷心菜和水果而非动物性食物为食。\par

49. 你以为树木的身体比我们的身体更好，其理由是否在于：肉体靠树木滋养，而树木却不靠肉体滋养？你忘记了明显的事实：植物在施以粪肥后，会变得更加茂盛肥沃，收成也更重；尽管你认为肉体是粪土的居所，这是你控诉肉体最严重的罪名。那么，这粪土却滋养了你视为洁净之物，尽管按你的说法，它是你所认为不洁之物中最不洁的部分。但如果你厌恶肉体，因为它源于性交，那么你应当对虫子的肉体感到满意——它们在水果中、在木头中、在地里大量繁殖，长得如此之大，且毫无性交。但这里有些不真诚。因为如果你因肉体由父母交合而成而不悦，你就不会说那些黑暗的王子是从他们自己树木的果实中出生的；因为你无疑对这些王子的评价比肉体更差，而肉体你尚且拒绝食用。\par

50. 你认为所有动物的灵魂都来自它们父母的食粮，从这束缚中你假装要解放那被捆绑在你食物中的神圣实体；这一想法与你不食肉类的做法完全矛盾，反而迫使你有责任去吃肉。因为如果灵魂是通过吃动物食物的人被束缚在身体里，你为何不抢先吃这食物以确保它们的解放？你回答说，那些人使之受束缚的善的部分并非来自动物食物，而是来自他们随肉一起食用的蔬菜。那么，对于只吃肉的狮子之灵魂，你将如何作答？你回答说，它们饮水，因而灵魂从水中被吸入并束缚在肉里。但数不清的鸟呢？只吃肉且无需饮水的鹰呢？在这里你束手无策，找不到答案。因为如果灵魂来自食物，而有些动物既不喝水，除肉外也无其他食物，却仍产下幼崽，那么肉中必有一定灵魂；你便当尝试通过吃肉来净化它的计划。或者你会说，猪因为吃蔬菜喝水，所以有光明之魂；而鹰因为只吃肉，所以有黑暗之魂，尽管它如此喜爱太阳？\par

51. 何等混乱的想法！何等惊人的愚昧！你若摒弃虚妄的虚构，而追随真理所准许的饮食节制，这一切原可避免；真理会教导你：应当避免奢侈的饮食，并非为避免玷污——因为并无此类事物——而是为克制感官的欲望。因为如果有人因不注意事物的本性以及灵魂与肉体的特性，而承认灵魂被肉食玷污，那么你也会承认灵魂远更被感官享受所玷污。那么，这是否合理？或者更确切地说，是否极其不合理：将一个也许为健康着想、毫无感官欲望而食用些许肉食的人，逐出精选者的行列；而若有人狼吞虎咽地食用加胡椒的块菌，你只能责备他过度，却不能谴责他滥用你的象征？因此，一个并非出于感官享受，而是为健康之故食用部分禽类的人，不能留在你的精选者中；而一个自愿放纵于无动物成分的蜜饯与糕点之过度食欲的人，却可以留下。你留住那深陷感官享受之污秽中的人，却驱逐那被你认为是仅因食物而玷污的人；尽管你承认感官享受的玷污远大于肉食的玷污。你紧抓住那个因美味蔬菜而喜悦得不能自持的人，却摒弃那个为充饥而取用任何拿到手且适于营养之物的人，不论食物还是排泄物他都准备使用。可赞的习俗！卓越的道德！著名的节制！\par

52. 再者，认为只有精选者才能接触作为食物摆上你餐桌、供你所谓净化之用的东西，这一观念导致了可耻的、有时甚至是罪恶的行为。因为有时所献之物如此之多，少数人难以轻易吃完；而将剩余之物给予他人，或至少将其丢弃，被视为亵渎，于是你不得不为了净化——如你所说——所有被献之物而饮食过量。然后，当你饱胀欲裂时，你残忍地强迫你教派的男孩们吃下剩余之物。因此，有人在罗马被指控强迫穷困儿童因这迷信理由而饮食致死。此事我不愿相信，若非我知道你们认为将这些食物给予非精选者，或至少将其丢弃，是何等有罪。那么唯一的办法就是吃掉它们；这每天导致贪食，有时还可能引致谋杀。\par

53. 基于同样理由，你们禁止给乞丐面包。为了显示怜悯，或者不如说是为了避免指责，你们建议给钱。这种残忍与愚蠢相等。因为假设有一个地方买不到食物：乞丐会饿死，而你们，以你们的智慧和仁慈，对黄瓜比对人更怜惜！这实在是（因为还能有什么更好的称呼）假装的怜悯，真正的残忍。然后注意其愚蠢。如果乞丐用你们给的钱自己买面包呢？你们所谓的\OriginalTerm{神圣部分}{divine part}，在他购买食物时，不也会遭受同样的痛苦吗，就像他从你们那里接受它作为礼物时所遭受的一样？所以这个有罪的乞丐将急于逃脱的天主的一部分拖入腐败，并借助你们的钱协助这一罪行！但你们以其大智大慧，认为只要不给即将谋杀的人一个可杀的人就够了，尽管你们明知给他钱去雇人被杀。还有什么疯狂能超过这个？结果是：如果他不能用钱得到食物，他就死了；如果他得到了食物，食物本身死了。一个是真正的谋杀；另一个是你们所谓的谋杀：虽然在两种情况下你们都犯了真正的谋杀罪。再者，允许你们的信徒吃动物肉，却禁止他们杀害动物，这是最大的愚蠢和荒谬。如果这种食物不污秽，你们自己吃吧。如果它污秽，还有什么比认为分离猪的灵魂与身体比用猪肉玷污人的灵魂更有罪更不合理呢？\par

% structure: p0073 source=h2 target=section reason=relative-heading-rank; base=h2

\section{第十七章——摩尼教徒中手的象征的描述}

54. 我们现在必须注意并讨论手的象征。首先，你们禁止杀害动物和伤害植物，这一点已被基督表明为纯粹的迷信；因为，基于我们与野兽和树木之间没有权利共同体，祂曾把魔鬼赶入一群猪中\BibleRef{玛 8:32}，并因一棵树上找不到果实而用咒诅使之枯萎\BibleRef{玛 21:19}。猪肯定没有犯罪，树也没有。我们并不疯狂到认为一棵树结不结果实是出于自己的选择。你们也不能回答说我们的主在这些行动中意在教导其他真理；因为人人都知道这一点。但可以肯定的是，天主之子不会为阐明真理而谋杀，如果你们把毁坏一棵树或一只动物称为谋杀的话。基督在人身上所行的奇迹——我们与人确实有权利共同体——是医治他们，而不是杀害他们。在野兽和树木身上也会一样，如果我们与它们有你们所想象的那种共同体的话。\par

55. 我认为在此引用圣经的权威是合适的，因为我们无法在此就动物的灵魂或树木中生命种类进行深入讨论。但既然你们保留着将圣经称为被篡改的权利，以防你们发现它过于强烈地反对你们——尽管你们从未断言关于树木和猪群的两段经文是伪经——但仍要提防有朝一日，当你们发现它们对你们多么不利时，你们可能也想对它们说同样的话。我将坚持我的计划，要求你们这些在证据和真理的提供上如此慷慨的人，首先告诉我，毁坏一棵树会造成什么伤害——我不是说摘一片叶子或一个苹果（因为如果你们中的一个人故意而不是意外地做了这件事，他立刻就会被判为滥用象征），而是说连根拔起一棵树。因为按照你们，树木中的灵魂，按你们的说法，是一个\OriginalTerm{理性灵魂}{rational soul}，当树被砍倒时，这个灵魂就从束缚中得到释放——在那里它遭受巨大痛苦却毫无收益。众所周知，在你们创始人的话语中，你们威胁将人变成树作为一种严厉的（尽管不是最严厉的）惩罚；并且树木中的灵魂不太可能像在人中那样增长智慧。不杀人的最大理由是，以免你们杀了一个其智慧或品德可能对许多人有用的人，或一个可能通过他人建议或自身心灵的神圣光照而获得智慧的人。灵魂离开身体时拥有的智慧越多，它的离去就越有益，正如我们从充分推理和广泛信仰中所知。因此，砍倒一棵树就是将一个灵魂从它无法在智慧上取得进展的身体中释放出来。你们——我指的是圣人——主要应该忙于砍树，并通过祈祷和圣咏将如此解放的灵魂引领到更好的境界。还是说，只有将灵魂吞进肚子里，而不是通过你们的理性来帮助它们，才能做到这一点？\par

56. 而且你们无法回避承认，树木中的灵魂在它们停留在那里期间在智慧上并无进展，当你们被问到为什么没有宗徒被派去教导树木，像教导人一样，或者为什么派给人间的宗徒也没有向树木传讲真理时，你们的回答必然是，这些灵魂在如此的身体中不能理解神圣诫命。但这个回答使你们陷入巨大困难；因为你们宣称这些灵魂能听见你们的声音并理解你们所说的话，能看见身体及其运动，甚至能洞察思想。如果这是真的，为什么它们不能从\OriginalTerm{光明宗徒}{apostle of light}那里学到任何东西？为什么它们甚至不能学得比我们更好，既然它们能看透心灵？你们的导师，按你们所说，他很难通过言语教导你们，却可能通过思想教导这些灵魂；因为它们能在他将思想表达出来之前就看到他心中的想法。但如果这是假的，你们就要想想自己陷入了怎样的错误。\par

57. 至于你们自己不摘水果或不拔蔬菜，却让你们的追随者去摘、去拔、并带给你们，以便你们不仅惠及带食物的人，也惠及所带的食物，有哪个有思想的人能忍受听这个？首先，不论你们自己犯下罪行，还是希望别人为你们犯下，这都无关紧要。\Quote{你们否认自己希望如此！}那么，如何能减轻莴苣和韭葱中所含的神圣部分的痛苦，除非有人拔下它们，带给圣人，让它们被净化？再者，如果你们经过一块田地，那里友谊的权利允许你们摘取任何想要的东西，你们看到一只乌鸦正要吃无花果，你们会怎么做？按照你们的想法，无花果本身岂不似乎在对你们说话，并可怜地恳求你们自己摘取它，把它埋葬在圣腹中，在那里它可以被净化并恢复过来，而不是让乌鸦吞下它，使它成为乌鸦被诅咒身体的一部分，然后它被移交到其他形态的束缚和折磨中？如果这是真的，你们多么残忍！如果不是，多么愚蠢！有什么比打破这个象征更与你们的意见相悖？有什么比保留它更不仁慈对待天主的肢体？\par

58. 这假定了你们虚假而空洞的想法是真实的。但你们可以被证明犯了明显而确实的残忍，这源于同一个错误。例如，假如有一个人躺在路上，身体因疾病而憔悴，因旅途而疲惫，因痛苦而半死，只能说出一些支离破碎的话，如果吃一个梨作为收敛剂对他有好处，他恳求你们在经过时帮助他，恳求你们从邻近的树上摘取那果子，没有神圣或世俗的禁令阻止你们这样做，而那个人肯定会因缺乏这果子而死，你们，一个基督徒和圣人，反而会继续前行，抛弃这样一个受苦哀求的人，唯恐那棵树会为失去它的果实而哀叹，而你们会注定遭受\Person{摩尼}因打破象征所威胁的惩罚。奇怪的风俗，奇怪的无害！\par

59. 现在，关于杀害动物以及你们意见的理由，此前所说的大部分也同样适用于此。因为杀死狼能对狼的灵魂造成什么伤害呢？既然狼活着就是狼，不会听任何宣讲者，也丝毫不会放弃流羊的血；而通过杀死它，理性灵魂，如你们所认为的，会从身体中的禁闭中得到释放？但你们甚至使你们的追随者不可以进行这种屠杀；因为你们认为这比砍树更糟糕。在这方面，你们的感官——即身体的感官——并没有太多可以指责的。因为我们通过它们的叫声看到和听到动物痛苦地死去，尽管人在禽兽身上忽略了这一点，因为禽兽没有理性灵魂，我们与它们没有权利的共融。但对于你们观察树木的感官，你们一定是完全瞎了。因为更且不说木头中没有表达痛苦的运动，还有什么比一棵树在翠绿茂盛、花团锦簇、果实累累时更好的呢？而这通常主要来自修剪。但如果它像你们所认为的那样感受到铁器，它应该因如此众多而严重的伤口而死，而不是在伤口处发芽，并以如此明显的喜悦复活。\par

60. 但为什么你们认为毁灭动物比毁灭植物罪更大，尽管你们认为植物比动物有更纯洁的灵魂？有一种补偿，我们被告知，当从田地中取来的一部分交给了选民和圣人进行净化时。这已经遭到驳斥；我认为已经充分证明，没有理由说蔬菜中比肉中有更多的善的成分。但是，如果有人以卖肉为生，并将生意的全部利润用于为你们的选民购买食物，并为那些圣人带来比任何农夫或农民都要多的供给，他难道不会用这个补偿作为他杀害动物的理由吗？但是，我们被告知，还有一些别的原因；因为狡猾的人在对无知的人讲话时，总能在自然的秘密中找到一些资源。那么，故事是这样的：那些从黑暗族类中被捉拿并捆绑的天上的王子，被世界的创造者分配在这个区域中一个地方，他们有特别属于他们的地上的动物，每个王子都有来自自己种类和阶级的动物；他们认为那些动物的屠杀者有罪，不允许他们离开大地，而是尽可能地用痛苦折磨来骚扰他们。哪个单纯的人不会被这个吓到，看见笼罩在这些事物上的黑暗中的虚无，而不会认为事实正如所描述的那样？但我要靠着天主的帮助，坚持我的目的，用最明显的真理来反驳神秘的虚假。\par

61. 那么，请问，如果陆地和海中的动物通过常规世代从这个王子族类中按普通繁衍方式延续，因为动物生命的起源可以追溯到那个族类的流产出生——请问，我再说，蜜蜂、青蛙以及许多其他不是通过性交产生的生物，是否可以杀害而免于惩罚？我们被告知不可以。所以你们禁止你们的追随者杀害动物，并不是因为它们与某些王子的关系。或者，如果你们对所有的身体都建立普遍关系，那么王子们同样会关心树木，而你们并不要求你们的追随者饶恕树木。你们又回到了那个软弱无力的答复，即对植物造成的伤害通过你们的追随者带到你们教会的水果来补偿。因为这意味着那些屠杀动物并在市场上卖肉的人，如果他们是你们的追随者，如果他们用所得的收入把蔬菜带给你们，就可以对每日的屠杀不以为意，并因你们的餐食而被免除其中可能有的任何罪。\par

62. 但若你说，为赎此杀戮之罪，那物必须作为食物给出，如水果和蔬菜那样——这是不可能的，因为选民不吃肉，所以你们的追随者不得宰杀动物——那么对于荆棘和杂草，农夫在清理田地时将其毁灭，却无法从中带给你任何食物，你又作何答复？这种丝毫不能滋养圣人的毁灭，如何能得赦免？也许你们也借着吃其中一些，为水果和蔬菜的好处，而消除所犯的罪。那么，若田地遭受蝗虫、家鼠或大鼠掠夺，正如我们常见的那样，你的乡下追随者能杀死它们而免于罪责，因为他为庄稼的利益而犯罪？此时你便茫然无措；因为你要么允许追随者宰杀动物——你的创始人禁止了这一点，要么禁止他们务农——而你的创始人使务农成为合法。事实上，你有时竟至于说，放高利贷者比农夫更为无害：你对瓜比对人类更怜惜。宁可伤害瓜，也不愿一个人负债破产。这是可向往、可称赞的正义，还是可憎恶、可诅咒的谬误？这是可称赞的怜悯，还是可憎恶的野蛮？\par

63. 再者，你们自己不戒除杀害虱、臭虫和跳蚤，又当如何？你们认为，说这些是我们身体的污秽，便是充分的借口。但对于跳蚤和臭虫，这显然不真；因为人人都知道这些动物并非来自我们的身体。此外，若你们如自己所宣称的那般憎恶性交，你们就该认为那些仅从我们身体生出而无其他生殖之动物更为洁净；因为它们虽生有自己的后代，却并非以平常的生殖方式由我们而生。再者，若我们必须认为活体的产生最为污秽，那么死体的产生必更为污秽。因此，杀死一只老鼠、蛇或蝎子——你们常说它们来自我们的死体——应少些害处。但且略过那些较不明显与确定的事，关于蜜蜂，有一种普遍看法是它们出自公牛尸体；故杀死它们也无害。或者，若此说亦被怀疑，至少人人都承认甲虫是由它们制作并掩埋的粪球中繁殖的。因此，你们应将这些动物以及其他不便详述的动物，视为比你们的虱更不洁；然而你们却认为杀死它们有罪，尽管不杀虱是愚蠢的。或许你们因虱之微小而轻视它们。但若动物以其大小被估量，你们就必须将骆驼置于人之上。\par

64. 此处我们可用那通阶层递进法——当我们还是你们的追随者时，此法常困扰我们。因为若一只跳蚤可因其微小而被杀，那么生在豆子中的苍蝇也可被如此对待。若此可，那么稍大一些的亦可，因其出生时甚至更小。然后，蜜蜂也可被杀，因其幼体不大于苍蝇。以此类推，直到蝗虫的幼体及蝗虫自身；再到老鼠的幼体及老鼠自身。简言之，显然我们最终可到达大象，故一个认为杀跳蚤因其微小而无罪的人，必得承认杀此庞然大物亦无罪。但我想，关于这些荒谬之处已说得足够。\par

% structure: p0085 source=h2 target=section reason=relative-heading-rank; base=h2

\section{第十八章——论胸部的象征及摩尼教徒的可耻奥秘。}

65. 最后，尚有胸部的象征，你们那大可怀疑的贞洁即在于此。因为你们虽不禁绝性交，却如宗徒早已说过的，禁止了本义的婚姻，尽管这是此类交合的唯一良好借口。无疑你们会对此抗议，并以我们不应指责你们——即你们高度赞赏并赞成完全的贞洁，却不禁止婚姻，因为你们的追随者（即你们当中的第二级者）被允许有妻子——来反责我们。在你们大声喧嚷、激动地说了这话之后，我将平静地问：岂不是你们主张生育子女——使灵魂被拘禁在肉身之中——比交合本身是更大的罪？岂不是你们常劝我们尽可能观察女子洁净后最易受孕的时期，并在那时戒除交合，以免灵魂被纠缠在肉身里？这证明你们赞成有妻子，不是为了生育子女，而是为了情欲的满足。在婚姻中，如婚姻法所宣布的，男女结合是为生育子女。因此，凡使生育子女成为比交合更大的罪者，便是禁止婚姻，并使女子不成为妻子，而成为情妇——为了一些馈赠而许给男子以满足其情欲。有妻子之处必有婚姻。但无生育之目的之处便无婚姻；故也无妻子。如此，你们禁止了婚姻。你们无法成功地辩护自己免于此控诉，这控诉早已由圣神预言般地加于你们。\par

66. 此外，当你们如此热切地渴望阻止灵魂通过夫妻交合而被束缚在肉身中，又如此热切地断言灵魂通过圣人的食物从种子中获得释放时，不幸的人们，你们难道不是默许了人们对你们的猜疑吗？\newline{}因为，关于谷粒、豆类、扁豆和其他种子，为何当你们吃它们时，你们想要释放灵魂，而对动物的种子却不这样？\newline{}因为你们对死兽之肉所说的——它是污秽的，因为其中没有灵魂——不能用于动物的种子；因为你们认为它紧紧拘禁着将在后代中出现的灵魂，并且你们承认\Person{摩尼}本人的灵魂就是这样被拘禁的。\newline{}既然你们的追随者无法将这些种子带给你们净化，谁会不怀疑你们秘密地在自己人当中进行这种净化，并向追随者隐瞒，以免他们离开你们呢？\newline{}即使你们不做这些事——但愿你们不做——你们仍然看到你们的迷信多么容易招致猜疑，而你们坚持认为通过饮食从身体和感官中释放灵魂，那么指责人们按照你们自己的主张去思想是多么不可能。\newline{}关于此事，我不想再多说什么：你们自己看到这里有多少可指责之处。但由于此事宁可被抑制也不应招致议论，又因为在我整个论述中，我的目的似乎是毫不夸张，只依据事实和论证，我们将转到其他事项。\par

% structure: p0088 source=h2 target=section reason=relative-heading-rank; base=h2

\section{第十九章——摩尼教徒的罪行。}

67. 那么，我们现在看到你们的三种象征的性质。这些就是你们的习俗。这就是你们那著名的规诫的结局，其中没有什么是确定的，没有什么坚定不移，没有什么前后一致，没有什么无可指摘，反而一切都是可疑的，或者更确切地说，毫无疑问是完全虚假的，一切都是自相矛盾、可憎、荒谬的。\newline{}总之，在你们的习俗中发现了如此众多且严重的恶行，以至于想要一一谴责它们的人，如果他能够稍作铺陈，至少需要为每一项分别撰写一篇专文。\newline{}如果你们遵守这些戒律并按你们的主张行事，那么没有任何幼稚、愚蠢或荒谬会超过你们；而当你们赞赏并教导这些事却不去实践时，你们所表现出的狡猾、欺骗和恶意，可以与任何能被描述或想象的事物相匹敌。\par

68. 在我极其认真和勤勉地追随你们的整整九年中，我没有听说你们的任何一位选民未违反这些规诫，或者至少未曾被怀疑如此。\newline{}许多人被发现在饮酒和吃肉，许多人去公共浴池；但这我们只是听说。\newline{}有些人被证实引诱他人妻子，以至于此事中我无法怀疑指控的真实性。\newline{}但假设这也只是传闻而非事实。\newline{}我自己亲眼看见——不仅是我，还有那些要么已经逃离那种迷信，要么我希望将要逃离的人——我们，我说，在\Place{迦太基}的一个广场上，在一条非常繁忙的道路上，不是一位，而是三位以上的选民走在我们后面，用如此猥亵的声音和手势与一些妇女搭讪，以至于超过了所有普通流氓的无耻和蛮横。\newline{}显然，这完全是他们的习惯，而且他们彼此之间也这样行事，因为无人因同伴在场而却步，这表明他们中大多数，如果不是全部，都染上了这种恶习。\newline{}因为他们并非同住一屋，而是住在相当不同的地方，并且完全是偶然地一起离开了他们相遇的地点。\newline{}这对我们打击很大，我们为此提出了投诉。\newline{}但谁曾想过施加惩罚——我不是说将他们逐出教会，甚至也不是按照罪行的严重程度予以严厉斥责？\par

69. 给予那些人免于惩罚的所有借口是：那时，当他们的聚会为法律所禁止时，担心受惩罚的人可能会通过告密进行报复。\newline{}那么，关于他们声称将永远在这个世界上遭受迫害——他们以此以为会得到更多尊重——又怎么说呢？\newline{}因为这就是他们对\BibleQuote{世界恨他们}这句话的应用。\BibleRef{若 15:18}\newline{}并且他们坚持认为真理必须在他们中间寻找，因为在关于圣神——护慰者的应许中，经上说\BibleQuote{世界不能接受祂}。\BibleRef{若 14:17}\newline{}这里不是讨论这个问题的地方。\newline{}但显然，如果你们将一直遭受迫害，直到世界终结，那么这种松懈以及所有这种不道德行为的无节制蔓延就不会有尽头，因为你们害怕得罪这种品性的人。\par

70. 当我们向最高当局报告一位妇女曾向我们投诉——她与其他妇女一起参加聚会，不怀疑这些人的圣洁——时，也给了我们同样的答复。\newline{}她投诉说，在聚会中，一些选民进来了，其中一人熄灭了灯，她无法辨认出的一个人试图拥抱她，若非她喊叫逃脱，那人就会强迫她犯罪。\newline{}我们必须推断，这种恶习多么普遍，以致导致了此事的恶行！\newline{}而这事是在你们守夜庆节的那一夜发生的。\newline{}事实上，除了害怕被告密之外，没有人能将犯罪者带到主教面前，因为他如此小心地避免被认出。\newline{}仿佛所有与他一起进来的人不都牵连在罪行中一样；因为在他们猥亵的欢乐中，他们都希望灯被熄灭。\par

71. 那时，当我们看到他们充满嫉妒、充满贪欲、渴望昂贵食物、不断争吵、对小事轻易激动时，何等大的怀疑之门被打开了！我们断定，如果他们找到隐秘或黑暗中的机会，他们就没有能力戒除他们所宣称要戒除的东西。有两个人品性相当好，思想活跃，是辩论中的领袖，我们与他们比其他人的关系更特别、更亲密。其中一人与我们交往甚密，因为他也从事自由技艺的学习；据说他现在是那里的长老。这两人彼此非常嫉妒，一人指控另一人——不是公开地，而是在谈话中，一有机会就耳语说——他曾对一位追随者的妻子进行犯罪攻击。而另一人在向我们辩解时，对另一位选民提出了同样的指控，这位选民与那位追随者同住，是他最信任的朋友。他突然闯进去，逮住了这人正与那女人在一起；而他的敌人和对手怂恿那女人和她的奸夫散布关于他的假报告，这样如果他提供任何信息就不会被相信。我们非常苦恼，深以为忧，尽管关于对那女人的攻击有疑问，但这两人的嫉妒之情——我们在那地方找不到比他们更好的人——表现得如此尖锐，不可避免地引起对其他事情的怀疑。\par

72. 另一件事是，我们经常在剧院里看到属于选民的人，有年纪的、被认为有品德的老人，还有一位白发苍苍的长老。我们略过那些年轻人，我们曾遇到他们为与舞台和赛马有关的人争吵；由此我们可以安全地推断，他们在私下里如何能够克制，因为他们无法抑制那种使他们暴露在追随者眼前、给他们带来耻辱和逃跑的激情。至于那位圣人，我们在无花果商贩街上参加他的讨论，如果他能使那位献身的贞女怀孕而不致使其受孕，他那令人发指的罪行会被发现吗？肿胀的子宫暴露了隐秘的、意想不到的邪恶。当她的兄弟，一个年轻人，从母亲那里听说这件事后，他深深感到伤害，但出于对宗教的尊重，克制了公开指控。他成功地将那人从那个教会中驱逐出去，因为这种行为不能总是被容忍；并且为了使罪行不至于完全不受惩罚，他安排了一些朋友将那人痛打一顿并踢打。当他受到这样的攻击时，他大声喊叫，要他们看\Person{摩尼}的意见的权威而饶恕他，即第一个英雄\Person{亚当}犯了罪，并在犯罪后成了更大的圣人。\par

73. 事实上，这就是你们关于\Person{亚当}和\Person{厄娃}的看法。说来话长；但我只触及与当前有关的事情。你们说，\Person{亚当}是从他的父母——流产的黑暗王子——产生的；他的灵魂里有着最多的光明部分，而很少对立种族。因此，当他过着圣洁的生活时，由于善的盛行，他里面的对立部分仍然被激起，以致他被引向夫妻交合。于是他堕落犯罪，但之后生活在更大的圣洁中。现在，我的抱怨不是主要关于这个邪恶的人，他在选民和圣人的外衣下，以令人震惊的不道德行为给一个陌生人家庭带来了这样的耻辱和谴责。我不以此责备你们。让这归因于那人的堕落品行，而非你们的习俗。我因那暴行责备那人，而非你们。然而，在你们所有人中，这一点在我看来是无法接受或容忍的：你们既主张灵魂是天主的一部分，却又坚持认为一点恶的混合胜过了更大力量和数量的善。谁相信这一点，当被激情煽动时，不会在这里找到借口，而不是克制和控制自己的激情？\par

% structure: p0096 source=h2 target=section reason=relative-heading-rank; base=h2

\section{第二十章．在罗马发现的丑行}

74. 关于你们的习俗，我还能再多说什么呢？我已经提到，当我在该城时，亲眼目睹了那些事情的发生。要详细叙述我不在时在罗马发生的一切，将耗费很长时间。不过，我会简短地讲述一下；因为那件事变得如此臭名昭著，以至于即使不在场的人也无法对其一无所知。后来我到了罗马，证实了我所听到的一切都是真实的，尽管这个故事是一位目击者告诉我的，我对他非常了解且极为敬重，以至于我丝毫不会怀疑。当时，你们的一位追随者，他完全配得上选民以他们闻名遐迩的克己生活，因为他受过良好的教育，并且热衷于捍卫你们的教派，他对有人当面指责那些生活在各种地方、四处投宿于最恶劣住所的选民的卑劣行为感到非常不满。因此，他渴望，如果可能的话，把所有愿意按照规诫生活的人聚集到自己的家中，并由他自己出资供养他们；因为他在花钱上的随意超乎常人，同时他拥有的可花费的财富也超乎常人。他抱怨说，他的努力受到了主教的疏忽的阻碍，他需要他们的协助才能成功。最后，终于找到了一位你们的主教——就我所知，此人非常粗鲁且不文雅，但不知怎的，正因其乖戾，反而更倾向于严格遵守道德规范。这位追随者热切地抓住此人，视之为他长久以来渴望并最终找到的人，并向他讲述了自己的全部计划。这位主教表示赞同并同意，率先搬进了那所房子。事情就这样做了，所有当时在罗马的选民都被聚集在那里。\WorkTitle{摩尼书信}中的生活规则被摆在他们面前。许多人认为难以忍受，便离开了；不少人感到羞愧，留了下来。他们开始按照约定的方式生活，正如这一最高权威所命令的。这位追随者一直热心地向每个人强调所有的事情，尽管在任何情况下，他本人从未承担过任何事情。与此同时，选民之间经常发生争吵。他们互相指控罪行，他听到这些感到痛心，并设法使他们在争吵中无意中相互揭露。揭露的内容卑劣得难以形容。这样，那些不同于其他人、愿意屈服于规诫的轭的人的真面目就显露出来了。那么，对于其他人，我们又该怀疑什么，或者更确切地说，该断定什么呢？最后，有一次他们聚集在一起，抱怨说他们无法遵守那些规定。随后便发生了反叛。这位追随者极其简明地陈述了他的立场：要么所有规定都必须遵守，要么那个认可这些无人能够履行的规诫的人一定是个大傻瓜。但是，正如不可避免的，一群人的狂乱叫喊压倒了一个人的意见。那位主教最终也屈服了，极其丢脸地逃走了；据说他违反规定偷偷弄到了食物，这些食物常常被发现。他用自己的私房钱提供了一笔资金，并小心地隐藏起来。\par

75. 如果你说这些事情是假的，那你就是在与过于明显和公开的事实相矛盾。但你如果愿意，尽可以这样说。因为，既然这些事情是确定的，并且愿意了解的人很容易就能知道，那么那些否认它们为真的人就显明了他们讲真话的习惯。不过，你们还有其他一些答复，我对此并不指责。因为你们要么说有些人确实遵守了你们的规诫，在谴责其他人时不应将他们与有罪的人混为一谈；要么说对你们教派成员品行的整个调查是错误的，因为问题在于该教义本身的品性。即使我同意这两点（尽管你们既无法指出那些忠诚遵守规诫的人，也无法使你们的异端摆脱所有这些轻浮和不义），我仍然必须坚持知道：为什么你们在看到一些人的不道德生活时，就责骂那些被称为天主教信徒的基督徒，而你们要么厚颜无耻地拒绝调查关于你们成员的情况，要么更厚颜无耻地不拒绝调查，希望人们理解为在你们稀少的成员中，有一些不知名的个体遵守了他们所宣称的规诫，但在天主教会的大批人群中却没有这样的人。\par

\SourceRecord{Translated by Richard Stothert.；Nicene and Post-Nicene Fathers, First Series；Vol. 4.；Edited by Philip Schaff.；Buffalo, NY: Christian Literature Publishing Co.,；1887.；<http://www.newadvent.org/fathers/1402.htm>.}
